% Part: many-valued-logic
% Chapter: three-valued-logics
% Section: goedel

\documentclass[../../../include/open-logic-section]{subfiles}

\begin{document}

\olfileid{mvl}{thr}{god}

\olsection{Gödel逻辑}

Kurt Gödel 引入了一系列 $n$ 值逻辑，其中每一种都包含直觉主义逻辑中的所有有效公式，并且包含于经典逻辑之中。第一个有趣的逻辑如下：

\begin{defn}\ollabel{defn:goedel}
\emph{三值Gödel逻辑}~$\LogGod$ 由如下矩阵定义：
\begin{enumerate}
  \item 标准命题语言 $\Lang L_0$，包含
  $\lfalse$、$\lnot$、$\land$、$\lor$、$\lif$。
  \item 真值集合 $V = \{\True, \Undef, \False\}$。
  \item $\True$ 是唯一的指定值，即 $V^+ = \{\True\}$。
  \item 对于 $\lfalse$，有 $\tf{\lfalse} = \False$。其余联结词的真值函数由下表给出：
  \begin{center}
    \begin{tabular}{c|c}
      $\tf{\lnot}[\LogGod]$ & \\
      \hline
      $\True$ & $\False$ \\
      $\Undef$ & $\False$ \\
      $\False$ & $\True$
    \end{tabular}
    \quad
    \begin{tabular}{c|ccc}
      $\tf{\land}[\LogGod]$ & $\True$ & $\Undef$ & $\False$ \\
      \hline
      $\True$ & $\True$ & $\Undef$ & $\False$ \\
      $\Undef$ & $\Undef$ & $\Undef$ & $\False$\\
      $\False$ & $\False$ & $\False$ & $\False$
    \end{tabular}
    \\[2ex]
    \begin{tabular}{c|ccc}
      $\tf{\lor}[\LogGod]$ & $\True$ & $\Undef$ & $\False$ \\
      \hline
      $\True$ & $\True$ & $\True$ & $\True$ \\
      $\Undef$ & $\True$ & $\Undef$ & $\Undef$ \\
      $\False$ & $\True$ & $\Undef$ & $\False$
    \end{tabular}
    \quad
    \begin{tabular}{c|ccc}
      $\tf{\lif}[\LogGod]$ & $\True$ & $\Undef$ & $\False$ \\
      \hline
      $\True$ & $\True$ & $\Undef$ & $\False$ \\
      $\Undef$ & $\True$ & $\True$ & $\False$  \\
      $\False$ & $\True$ & $\True$ & $\True$
    \end{tabular}
  \end{center}
\end{enumerate}
\end{defn}

你会注意到，$\land$ 和~$\lor$ 的真值表与 \L ukasiewicz 逻辑和强 Kleene 逻辑中的相同，但 $\lnot$ 和~$\lif$ 的真值表在每种逻辑中各不相同。在Gödel逻辑中，$\tf{\lnot}(\Undef) = \False$。与 \L ukasiewicz 逻辑和 Kleene 逻辑不同，$\tf{\lif}(\Undef, \False) = \False$；与 Kleene 逻辑不同（但与 \L ukasiewicz 逻辑相同），$\tf{\lif}(\Undef,
\Undef) = \True$。

如前文所述其与直觉主义逻辑的联系所暗示的，$\LogGod[3]$ 接近直觉主义逻辑。所有直觉主义有效式都是~$\LogGod[3]$ 中的重言式，并且许多在直觉主义中无效的经典重言式在~$\LogGod[3]$ 中也不是重言式。例如，以下公式不是重言式：
\begin{align*}
  & p \lor \lnot p && (p \lif q) \lif (\lnot p \lor q) \\
  & \lnot\lnot p \lif p && \lnot(\lnot p \land \lnot q) \lif (p \lor q) \\
  & ((p \lif q) \lif p) \lif p && \lnot(p \lif q) \lif (p \land \lnot q)
\end{align*}
然而，并非 $\LogGod[3]$ 中的所有重言式都是直觉主义有效的，例如 $\lnot\lnot p \lor \lnot p$ 或 $(p
\lif q) \lor (q \lif p)$。

\begin{prob}
  给出真值表以证明下列公式是~$\LogGod[3]$ 的重言式：
  \begin{align*}
    & \lnot\lnot p \lor \lnot p\\
    & (p \lif q) \lor (q \lif p) \\
    & \lnot(p \land q) \lif (\lnot p \lor \lnot q) \\
    & (p \lif q) \lor (q \lif r) \lor (r \lif s)
  \end{align*}
\end{prob}

\begin{prob}
  给出真值表以证明下列公式不是~$\LogGod[3]$ 的重言式：
  \begin{align*}
    & (p \lif q) \lif (\lnot p \lor q) \\
    & \lnot(\lnot p \land \lnot q) \lif (p \lor q) \\
    & ((p \lif q) \lif p) \lif p \\
    & \lnot(p \lif q) \lif (p \land \lnot q)
  \end{align*}
\end{prob}

\begin{prob}
  下列语义后承关系在Gödel逻辑中哪些成立？对每一项给出真值表。
  \begin{enumerate}
    \item $p, p \lif q \Entails q$
    \item $p \lor q, \lnot p \Entails q$
    \item $p \land q \Entails p$
    \item $p \Entails p \land p$
    \item $p \Entails p \lor q$
  \end{enumerate}
\end{prob}

\end{document}